% SOURCE_STATEMENT: "Determine the connection between Mock modular forms and geometric invariants like Gromov-Witten, Donaldson, and OSV conjecture."
% Preparation and final self-audit: CONTRIBUTING.md prompts C and D; source/audit artifacts retained in the local physics_statements_2026_10_11 batch.
% !TeX program = lualatex
\documentclass[11pt]{article}
\usepackage[margin=25mm]{geometry}
\usepackage{fontspec}
\setmainfont{Latin Modern Roman}
\usepackage{amsmath,amssymb,mathtools}
\usepackage{unicode-math}
\setmathfont{Latin Modern Math}
\usepackage{xurl,hyperref}
\hypersetup{colorlinks=true,urlcolor=blue}
\setcounter{secnumdepth}{0}
\setlength{\parindent}{0pt}
\setlength{\parskip}{5pt}
\setlength{\emergencystretch}{3em}
\newcommand{\sref}[2]{\hyperlink{source:#1:#2}{[S#2]}}
\newcommand{\cataloguescope}[1]{\paragraph{Recorded target.}#1}
\newcommand{\R}{\mathbb R}
\newcommand{\C}{\mathbb C}
\newcommand{\N}{\mathbb N}
\newcommand{\Z}{\mathbb Z}
\newcommand{\Q}{\mathbb Q}
\newcommand{\F}{\mathbb F}
\newcommand{\eps}{\varepsilon}
\begin{document}
\section{Mock modularity, geometric invariants, and finite pole tomography}\label{statement}

\subsection{Definitions and mathematical statement}

\paragraph{Shared notation.}$\mathbb N=\{1,2,\ldots\}$; the standard number fields have their usual meanings. All additional conventions are stated locally.


\paragraph{Definitions and notation.}Gromov--Witten invariants are virtual counts of stable maps into a specified target. Donaldson-type invariants depend on the relevant geometric/gauge-theoretic setting. Mock modular forms are holomorphic parts of modular completions. The OSV conjecture relates a black-hole counting ensemble to topological-string data under particular physical and normalization assumptions.

\paragraph{Catalogue wording retained for traceability.}
Determine the connection between Mock modular forms and geometric invariants like Gromov-Witten, Donaldson, and OSV conjecture.
\paragraph{Scope and source qualification.}The source uses Donaldson without specifying the variant; Donaldson, Donaldson–Thomas and Gromov–Witten invariants are not interchangeable. The record remains a broad connection programme and does not assert an unspecified universal equivalence. The catalogue title/status was paired with a machine-generated research attempt; that attempt is not peer-reviewed evidence of settlement. The statement covers the original source questions rather than only the title’s subcase.
\subsection{Short English statement}
Explain links between mock modular forms, geometric counting invariants and the OSV framework.
\subsection{Sources}
\begin{itemize}
\item[\hypertarget{source:statement:1}{S1}] ULAM.AI and the UnsolvedMath Contributors, UnsolvedMath, catalogue entry. Catalogue curation under CC BY 4.0; definitions and English exposition added locally.
\url{https://huggingface.co/datasets/ulamai/UnsolvedMath}.
\item[\hypertarget{source:statement:2}{S2}] American Institute of Mathematics, Gromov-Witten invariants and number theory, source item 5.3.
\url{http://aimpl.org/gromwitnumthry/5/}.
\end{itemize}
\label{end:statement}
\end{document}
